\section*{Chapitre \ref{ch:b2:meiosis-heredity} --- Méiose, brassage génétique et hérédité}

\begin{solution}{exo:b2:meiosis-heredity:1}
Mitose~: une division, pas d'appariement des homologues, deux cellules,
chacune génétiquement identique à la cellule mère ($2n \to 2n$).
\omterm{def:b2:meiosis-heredity:meiosis}{Méiose}~: deux divisions après une seule réplication, appariement et
crossing-over des homologues, quatre cellules, \omterm{def:b2:meiosis-heredity:meiosis}{haploïdes} et toutes
génétiquement différentes ($2n \to n$).
\end{solution}

\begin{solution}{exo:b2:meiosis-heredity:2}
$2^{23} = 8.4\times 10^{6}$~; $2^{4} = 16$~; $2^{7} = 128$ --- avant
que le crossing-over ne multiplie chacun de ces nombres sans limite.
\end{solution}

\begin{solution}{exo:b2:meiosis-heredity:3}
$Tt\times Tt$~: \omterm{def:b2:meiosis-heredity:vocabulary}{génotypes} $1\,TT : 2\,Tt : 1\,tt$, \omterm{def:b2:meiosis-heredity:vocabulary}{phénotypes} $3$ hautes
pour $1$ naine. $Tt\times tt$~: $1\,Tt : 1\,tt$, soit $1 : 1$. Croiser
la plante haute avec une plante naine~: une descendance entièrement
haute signifie $TT$, une moitié de naines signifie $Tt$
(ou l'autoféconder~: $TT$ donne une descendance homogène, $Tt$ ségrège
en $3 : 1$).
\end{solution}

\begin{solution}{exo:b2:meiosis-heredity:4}
Des gènes liés sont portés par le même chromosome, et leurs combinaisons
alléliques parentales sont surreprésentées parmi les gamètes~; la
\omterm{def:b2:meiosis-heredity:linkage}{fréquence de recombinaison} est la proportion de gamètes recombinés
(descendants recombinés d'un croisement-test)~; un \omterm{def:b2:meiosis-heredity:linkage}{centimorgan}
correspond à un pour cent de recombinaison. Des gènes éloignés sur un
long chromosome subissent au moins un crossing-over entre eux dans
presque chaque \omterm{def:b2:meiosis-heredity:meiosis}{méiose}, et les crossing-over multiples rendent les
combinaisons aléatoires~: $r$ atteint alors $\tfrac12$ ---
indiscernable d'un brassage indépendant.
\end{solution}

\begin{solution}{exo:b2:meiosis-heredity:5}
Total 556~; effectifs attendus $312.75 : 104.25 : 104.25 : 34.75$~; $\chi^2 =
2.25^2/312.75 + 3.75^2/104.25 + 3.25^2/104.25 + 2.75^2/34.75 = 0.016
+ 0.135 + 0.101 + 0.218 = 0.47 < 7.81$~: compatible avec $9 : 3 : 3 :
1$.
\end{solution}

\begin{solution}{exo:b2:meiosis-heredity:6}
Contre $1 : 1 : 1 : 1$ (250 chacune)~: $\chi^2 = (162^2 + 138^2 + 147^2
+ 153^2)/250 = 361 \gg 7.81$~: les gènes sont liés. Recombinés $103 + 97 = 200$
sur 1000~: $r = 0.20$, soit \qty{20}{cM}. En répulsion, $Ab/aB$ donne 400
$Ab$, 400 $aB$, 100 $AB$, 100 $ab$~: la même distance, les classes
parentales et recombinées étant échangées.
\end{solution}

\begin{solution}{exo:b2:meiosis-heredity:7}
$d = -\tfrac12\ln(1 - 2r)$~: $r = 0.20$ donne $-\tfrac12\ln 0.6 =
0.255$, soit \qty{25.5}{cM}~; $r = 0.45$ donne $-\tfrac12\ln 0.1 = 1.15$,
soit \qty{115}{cM}. $r = \tfrac12(1 - e^{-2d})$~: \qty{30}{cM} donne
$\tfrac12(1 - e^{-0.6}) = 0.226$~; \qty{100}{cM} donne $\tfrac12(1 -
e^{-2}) = 0.43$. En dessous de \qty{10}{cM}, la correction est
négligeable~; au-delà de \qty{30}{cM}, elle est indispensable, et
au-delà de \qty{50}{cM}, $r$ sature, si bien qu'un simple croisement à
deux points ne peut plus mesurer la distance --- on cartographie les
gènes éloignés en enchaînant des gènes proches.
\end{solution}

\begin{solution}{exo:b2:meiosis-heredity:8}
(a) $X^wX^w \times X^+Y$~: tous les fils blancs ($X^wY$), toutes les
filles rouges ($X^+X^w$). (b) $X^+X^w \times X^wY$~: fils pour moitié
rouges, pour moitié blancs~; filles pour moitié rouges ($X^+X^w$), pour
moitié blanches ($X^wX^w$). (c) filles
$X^+X^w$ $\times$ frères $X^wY$~: comme en (b) --- dans chaque sexe,
moitié blancs, moitié rouges.
\end{solution}

\begin{solution}{exo:b2:meiosis-heredity:9}
Deux enzymes en série fabriquent le pigment~; la couleur pourpre exige
un \omterm{def:b2:meiosis-heredity:vocabulary}{allèle} \omterm{def:b2:meiosis-heredity:vocabulary}{dominant} aux deux loci ($C\_P\_$). Les lignées sont $CCpp$ et
$ccPP$ (chacune blanche pour une raison différente)~; l'hybride $CcPp$
est pourpre~; autofécondé, il donne $9\,C\_P\_$ pourpres pour $7$
blanches ($3\,C\_pp + 3\,ccP\_ +
1\,ccpp$). Croisement-test $CcPp\times ccpp$~: $1$ pourpre ($CcPp$) pour
$3$ blanches.
\end{solution}

\begin{solution}{exo:b2:meiosis-heredity:10}
Parentaux $+++$, $abc$~; doubles crossing-over $+b+$, $a+c$~; $b$ est le
gène qui a changé~: ordre $a$--$b$--$c$. $r_{ab} = (11 + 9 + 2)/1202
= 0.018$ (\qty{1.8}{cM})~; $r_{bc} = (118 + 112 + 2)/1202 = 0.193$
(\qty{19.3}{cM}). Doubles attendus $0.018\times 0.193\times 1202 =
4.2$~; observés 2~: coïncidence $0.47$, interférence $0.53$.
\end{solution}

\begin{solution}{exo:b2:meiosis-heredity:11}
Asques de seconde division~: $300/1000 = \qty{30}{\%}$~; distance $=
\tfrac12\times 30 = \qty{15}{cM}$. Un crossing-over entre le gène et le
centromère n'implique que deux des quatre chromatides~: un asque à
ségrégation de seconde division porte donc deux chromatides
recombinées et deux parentales, et la \omterm{def:b2:meiosis-heredity:linkage}{fréquence de recombinaison} vaut
la moitié de la fréquence de ces asques. Pour $2 : 4 : 2$~: un
crossing-over entre le locus et le centromère échange les deux
chromatides internes, si bien que chaque demi-fuseau de la \omterm{def:b2:meiosis-heredity:meiosis}{méiose} I
porte une chromatide $A$ et une chromatide $a$~; la \omterm{def:b2:meiosis-heredity:meiosis}{méiose} II les
sépare ensuite dans l'ordre $A\,a \mid
a\,A$ (ou l'inverse), et la mitose finale double chaque spore~:
$2 : 4 : 2$.
\end{solution}

\begin{solution}{exo:b2:meiosis-heredity:12}
Sa mère est obligatoirement conductrice (elle a reçu l'unique X de son
père)~: la femme a donc reçu l'\omterm{def:b2:meiosis-heredity:vocabulary}{allèle} avec une probabilité $\tfrac12$.
Deux fils en bonne santé~: une conductrice a une probabilité
$(\tfrac12)^2 = \tfrac14$ d'avoir deux fils sains, une non-conductrice
une probabilité 1. Formule de Bayes~:
$P(\text{conductrice}\mid\text{données}) = (\tfrac12\times\tfrac14)/(\tfrac12
\times\tfrac14 + \tfrac12\times 1) = \tfrac{1/8}{5/8} = \tfrac15$.
\end{solution}

\begin{solution}{pb:b2:meiosis-heredity:1}
\textbf{1.} Parents $vg/vg\;+/+$ et $+/+\;e/e$~; $F_1$ $+/vg\;+/e$~;
gamètes $+\,+$, $+\,e$, $vg\,+$, $vg\,e$, chacun $\tfrac14$.
\textbf{2.} $900 : 300 : 300 : 100$.
\textbf{3.} $\chi^2 = 8^2/900 + 9^2/300 + 4^2/300 + 3^2/100 = 0.07 +
0.27 + 0.05 + 0.09 = 0.48 < 7.81$~: le brassage indépendant est
confirmé.
\textbf{4.} Les \omterm{def:b2:meiosis-heredity:vocabulary}{homozygotes} sauvages aux deux loci représentent
$\tfrac{1}{16}$ du total, et $\tfrac{9}{16}$ sont sauvages~: $\tfrac19$.
\textbf{5.} Sauvages, vestigiales, ébène, vestigiales ébène, $\tfrac14$
chacune.
\textbf{6.} Le \omterm{def:b2:meiosis-heredity:vocabulary}{phénotype} résulte du \omterm{def:b2:meiosis-heredity:vocabulary}{génotype} et de l'environnement (la
température agit sur les enzymes du pigment)~: un \omterm{def:b2:meiosis-heredity:vocabulary}{génotype} définit une
norme de réaction, non un caractère figé.
\textbf{7.} $b\,+/+\,vg$ (répulsion)~: $b$ est venu avec $+_{vg}$ d'un
parent, $vg$ avec $+_b$ de l'autre.
\textbf{8.} Parentaux~: noires ($b\,+$) 965 et vestigiales ($+\,vg$)
944~; recombinés~: noires vestigiales ($b\,vg$) 206 et sauvages
($+\,+$) 185. Le chromosome sauvage $+\,+$ n'existait pas chez la
$F_1$~; il ne peut être produit que par un crossing-over.
\textbf{9.} $r = 391/2300 = 0.17$~: \qty{17}{cM}.
\textbf{10.} $d = -\tfrac12\ln(1 - 0.34) = -\tfrac12\ln 0.66 = 0.21$~:
\qty{21}{cM}.
\textbf{11.} Il n'y a pas de crossing-over chez le mâle de
\emph{Drosophila}~: les gamètes du mâle ne portent que les deux
chromosomes parentaux.
\textbf{12.} $F_1$ $b\,vg/+\,+$ (couplage)~: classes parentales sauvage
et noire vestigiale, \qty{41.5}{\%} chacune~; classes recombinées noire
et vestigiale, \qty{8.5}{\%} chacune.
\textbf{13.} $X^hX^+$~: une fille reçoit l'unique X de son père, qui
porte $h$, et un X normal de sa mère non atteinte.
\textbf{14.} $\tfrac12$ (fils) $\times$ $\tfrac12$ (reçoit $X^h$) $=
\tfrac14$.
\textbf{15.} $\tfrac12$~: elle reçoit l'un ou l'autre des X de sa
mère~; le père donne un X normal.
\textbf{16.} Chaque enfant est non atteint avec une probabilité
$\tfrac34$~: $(\tfrac34)^3 = \tfrac{27}{64} = 0.42$.
\textbf{17.} Un fils reçoit le Y de son père, jamais son X. Une femme
n'est atteinte qu'à l'état $X^hX^h$~: il faut un père atteint et une
mère conductrice (ou atteinte), ce qui est rare.
\textbf{18.} Probabilité a priori $\tfrac12$. Deux fils sains ont une
probabilité $\tfrac14$ si elle est conductrice, 1 sinon~: probabilité
a posteriori $(\tfrac12\times
\tfrac14)/(\tfrac12\times\tfrac14 + \tfrac12) = \tfrac15$.
\textbf{19.} L'unique X d'un fils vient de sa mère, et le Y ne porte
aucun de ces gènes~: chaque fils exprime sans masque un gamète maternel,
et le croisement est son propre croisement-test.
\textbf{20.} Parentaux $+++$ (853), $ywm$ (833)~; doubles crossing-over
$+w+$ (2), $y+m$ (2)~; $w$ a changé~: ordre $y$--$w$--$m$.
\textbf{21.} $r_{yw} = (24 + 26 + 2 + 2)/2000 = 0.027$~; $r_{wm} = (128
+ 132 + 2 + 2)/2000 = 0.132$. Carte~: $y$ --- \qty{2.7}{cM} --- $w$ ---
\qty{13.2}{cM} --- $m$.
\textbf{22.} Attendus $0.027\times 0.132\times 2000 = 7.1$~; observés
4~: coïncidence $0.56$, interférence $0.44$.
\textbf{23.} Recombinés $y$--$m$~: $y++$, $+wm$, $yw+$, $++m$~: $(24 +
26 + 128 + 132)/2000 = 0.155$, contre $0.027 + 0.132 = 0.159$~: les
quatre doubles crossing-over restaurent la combinaison parentale
$y$--$m$ et échappent au comptage, soit $2\times 4/2000 = 0.004$.
\textbf{24.} Les filles reçoivent l'X $y\,w\,m$ du père et un X
maternel~: toutes sont \omterm{def:b2:meiosis-heredity:vocabulary}{hétérozygotes} et de \omterm{def:b2:meiosis-heredity:vocabulary}{phénotype} sauvage, quel que
soit l'X recombiné qu'elles portent~; la recombinaison est invisible.
\textbf{25.} $y$ --- \qty{2.7}{cM} --- $w$ --- \qty{13.2}{cM} --- $m$~;
coefficient de coïncidence $0.56$, interférence $0.44$.
\end{solution}
